\chapter{Requisitos detallados y validación con cliente}
\label{sec:anexo_requisitos_detallados}

Este anexo recoge de forma más detallada la información utilizada para construir el
registro de requisitos del proyecto y el proceso de validación seguido con el
cliente. Su objetivo es dejar constancia del origen de las decisiones de diseño y
de la evolución del alcance tras contrastar la propuesta inicial.

\section{Resumen de reuniones con el cliente}

\begin{table}[H]
  \centering
  \caption{Resumen de reuniones mantenidas con el cliente.}
  \label{tab:reuniones_cliente}
  \begin{tabularx}{\textwidth}{p{3.3cm}X X}
    \toprule
    \textbf{Reunión} & \textbf{Objetivo} & \textbf{Resultado principal} \\
    \midrule
    Toma inicial de requisitos &
    Comprender el problema, el alcance esperado y las limitaciones de la herramienta
    existente. &
    Se identificó la necesidad de una solución formada por aplicación móvil,
    portal web, funcionamiento offline, grupos de usuarios, mapas, dashboards y
    taxonomía de residuos. \\[4pt]

    Revisión de requisitos y mockups &
    Contrastar los requisitos definidos con los prototipos interactivos del portal
    y la aplicación móvil. &
    Se validó que los mockups cubrían los procesos principales y se reforzó la
    prioridad de la captura offline, la separación por grupos y la trazabilidad de
    los datos científicos. \\
    \bottomrule
  \end{tabularx}
\end{table}

\section{Requisitos obtenidos en la reunión inicial}

Durante la primera reunión se extrajeron los siguientes requisitos de alto nivel:

\begin{enumerate}
  \item Desarrollar una plataforma compuesta por aplicación móvil y portal web.
  \item Permitir registrar residuos en distintos entornos naturales, no solo ríos y
    mares, sino también playas, bosques y otros terrenos.
  \item Registrar información científica mínima de cada observación: tipo de
    residuo, cantidad, posición, fecha, hora, entorno y comentarios.
  \item Permitir adjuntar fotografías a las observaciones.
  \item Garantizar que la aplicación móvil pueda utilizarse sin conexión a internet.
  \item Sincronizar automáticamente los datos cuando el dispositivo recupere
    conectividad.
  \item Gestionar usuarios, grupos y permisos para mantener la privacidad de los
    datos de cada equipo.
  \item Ofrecer en el portal web herramientas de consulta, mapas, dashboards y
    exportación de datos.
  \item Mantener una taxonomía de residuos jerárquica, navegable y ampliable.
  \item Preservar trazabilidad científica mediante versiones de taxonomía, metadatos
    de sesión y registros de auditoría.
\end{enumerate}

\section{Validación de requisitos mediante mockups}

La segunda reunión se centró en revisar los prototipos funcionales y comprobar que
los flujos implementados representaban correctamente los requisitos ya recogidos.
En lugar de generar un nuevo bloque de requisitos, esta revisión permitió confirmar
la cobertura de los casos de uso principales:

\begin{itemize}
  \item Inicio de sesión y diferenciación entre perfiles de administrador y
    colaborador.
  \item Selección de grupo activo y filtrado de la información asociada.
  \item Creación de sesiones de muestreo en la aplicación móvil.
  \item Registro de observaciones con taxonomía, cantidad, tamaño, color,
    comentarios, GPS y fotografía.
  \item Visualización de observaciones y sesiones desde el portal web.
  \item Consulta geográfica mediante mapas.
  \item Representación del estado de sincronización y de los datos pendientes.
\end{itemize}

Los ajustes visuales e internos realizados durante la construcción de los mockups
se consideran decisiones de diseño del equipo y no requisitos adicionales
solicitados por el cliente. Por ello no se documentan como cambios de alcance, sino
como parte del proceso de refinamiento del prototipo.

\section{Requisitos funcionales detallados}

El registro de requisitos funcionales queda agrupado en seis bloques:

\begin{itemize}
  \item \textbf{Autenticación y sesión:} inicio de sesión, cierre de sesión y
    recuperación de contraseña.
  \item \textbf{Roles, grupos y permisos:} gestión de grupos, usuarios, membresías y
    selección de grupo activo.
  \item \textbf{Captura de campo:} sesiones de muestreo, observaciones, GPS,
    fotografías, cantidad, tamaño, color y comentarios.
  \item \textbf{Taxonomía:} navegación jerárquica, búsqueda y favoritos.
  \item \textbf{Sincronización:} funcionamiento offline, cola de sincronización,
    sincronización manual y automática, y estados por elemento.
  \item \textbf{Portal web:} listados, filtros, mapas, heatmaps, dashboards,
    fotografías, exportación, grupos, usuarios y configuración.
\end{itemize}

\section{Requisitos no funcionales detallados}

Los requisitos no funcionales se agrupan en:

\begin{itemize}
  \item \textbf{Seguridad:} comunicación cifrada, control de acceso por grupo y
    protección de adjuntos.
  \item \textbf{Fiabilidad:} sincronización reintentable y preservación de datos ante
    fallos o cierres inesperados.
  \item \textbf{Usabilidad:} interfaz adaptada a uso de campo en móvil y tablet.
  \item \textbf{Operación:} logs, copias de seguridad y matriz requisito-prueba.
  \item \textbf{Compatibilidad:} soporte para dispositivos móviles modernos, con
    especial atención a Android e iOS.
\end{itemize}

\section{Requisitos de datos científicos}

Los requisitos de datos científicos se centran en garantizar que las observaciones
registradas puedan analizarse posteriormente sin perder contexto:

\begin{itemize}
  \item Taxonomía jerárquica con códigos estables.
  \item Metadatos de muestreo por sesión.
  \item Metadatos de foto cuando estén disponibles.
  \item Extensión versionada de la taxonomía.
  \item Validaciones de obligatoriedad, rango y consistencia.
  \item Exportaciones con diccionario de columnas y unidades.
  \item Versión de taxonomía y de aplicación asociada a sesiones/exportaciones.
\end{itemize}
